\chapter{Thích Mà Không Nói Ra}
\chapterintro{Đỏ mặt nhưng vẫn giả vờ}{Nói ra thì ngại, im luôn thì nhớ. Chương này là những lần thương thầm rất học trò: viết nháp tên người ta, chờ ở hành lang, nhắn ngủ sớm mà tim tự loạn.}

\poem{Bài 16}{Bài 16 - Tên ai kín giấy nháp}{A secret name fills the draft page before the lesson notes do.}{Thích mà không nói nhé\\Nên đành viết ở giấy nháp\\Bài hôm nay còn trống nửa\\Tên ai thì kín cả trang}

\poem{Bài 17}{Bài 17 - Đợi ở cuối hành lang}{Waiting at the corridor corner is a full ritual of teenage courage.}{Thích mà không nói nhé\\Nên hay đứng cuối hành lang\\Đợi ai đi ngang một chút\\Rồi làm như thể tình cờ}

\poem{Bài 18}{Bài 18 - Nhắn ngủ sớm thôi nha}{One simple bedtime text can keep someone awake for hours.}{Thích mà không nói nhé\\Chỉ dám nhắn ngủ sớm nha\\Ai ngờ mình vui quá mức\\Thức luôn đến tận mười hai}

\poem{Bài 19}{Bài 19 - Ảnh lớp mà nhìn một người}{In a class photo, the heart never scans the whole frame.}{Thích mà không nói nhé\\Chụp ảnh lớp đứng rất ngoan\\Mà về xem đi xem lại\\Toàn phóng to đúng một người}

\poem{Bài 20}{Bài 20 - Muốn đi chung một ô}{Half an umbrella is enough to start a whole season of remembering.}{Thích mà không nói nhé\\Nên trời mưa thấy vui hơn\\Chỉ mong ai kia lên tiếng\\Đi chung cho đỡ ướt đường}
